\documentclass[11pt,a4paper]{article}
\usepackage{fontspec}
\usepackage[french]{babel}
\usepackage[margin=2.2cm]{geometry}
\usepackage{amsmath}
\usepackage{xcolor}
\setlength{\parskip}{0.5em}
\setlength{\parindent}{0pt}

\newcommand{\partie}[3]{\bigskip\noindent\colorbox{gray!15}{\parbox{\dimexpr\linewidth-2\fboxsep}{%
  \textbf{#1}\hfill\textbf{#2}\\ \small\textit{À l'écran :} #3}}\par}

\title{Script de la vidéo (10 minutes maximum)}
\author{Projet IMA -- édition de Poisson par Fourier}
\date{}

\begin{document}
\maketitle

\textbf{Principe.} Neuf parties, environ 9 min 20 au total, pour garder une marge sous les 10 minutes.
Le texte est écrit pour être dit à voix haute, à un rythme posé (environ 140 mots par minute) ; il n'a
pas besoin d'être lu mot pour mot. À l'écran : les figures du rapport (dossier \texttt{docs/rapport/figures/}), en plein
écran ou deux par deux. Le temps indiqué à droite est le moment où la partie commence.

%------------------------------------------------------------------
\partie{1. Le problème}{0:00 -- 0:50}{titre du projet, puis la figure 9 de l'article et mon résultat
(\texttt{article\_collage.png}, colonne de droite).}

Bonjour. Mon projet porte sur l'article de Morel, Petro et Sbert, « Fourier implementation of Poisson
image editing ». L'édition de Poisson consiste à modifier une image non pas en changeant ses valeurs, mais
en changeant son gradient : on impose le gradient voulu, et on cherche l'image dont le gradient en est le
plus proche. C'est ce qui permet de coller un objet sans couture, ou d'éclaircir une zone sans créer de
bord. L'idée de l'article est de résoudre ce problème sur toute l'image, d'un seul coup, par transformée
de Fourier. J'ai réimplémenté la méthode, reproduit ses expériences, analysé ses défauts, et je l'ai
appliquée à des images médicales, en 2D et en 3D.

%------------------------------------------------------------------
\partie{2. La méthode}{0:50 -- 2:00}{l'équation $\min\int|\nabla u - V|^2$, puis $\Delta u =
\operatorname{div}V$, puis l'éq.~12 encadrée. Éventuellement un schéma : image, image miroir $2H\times2W$.}

On cherche l'image $u$ qui minimise l'écart entre son gradient et un champ $V$. Cela revient à résoudre
l'équation de Poisson : le laplacien de $u$ égale la divergence de $V$, avec une condition de Neumann au
bord.

Pour la résoudre par Fourier, on prolonge l'image par symétrie miroir : on obtient une image deux fois
plus grande, périodique et sans saut, ce qui donne la condition de Neumann gratuitement. En Fourier,
dériver revient à multiplier par $iw$. L'équation de Poisson devient donc une simple division, coefficient
par coefficient : c'est l'équation 12. Trois FFT, un coût en $N\log N$, quelle que soit la forme de la
zone modifiée.

En codant, j'ai dû régler quatre points que l'article ne détaille pas : le recentrage des fréquences, un
signe, la fréquence nulle, et le retour du résultat dans $[0,255]$.

%------------------------------------------------------------------
\partie{3. Vérification}{2:00 -- 2:40}{le tableau 1 du rapport (MSE, PSNR).}

Premier test : si on donne au solveur le gradient de l'image elle-même, on doit retrouver l'image. C'est
le cas à la précision de la machine : un PSNR d'environ 320 décibels, contre 70 à 96 dans l'article. Et
si on remplace le gradient spectral par une simple différence finie, on tombe à 39 décibels, avec des
erreurs de plus de 100 niveaux de gris sur les contours. La raison : la différence finie mesure la pente
entre deux pixels, et le solveur, qui divise par $iw$, reconstruit alors une image décalée d'un
demi-pixel et floutée. Le gradient et le solveur doivent être exactement l'inverse l'un de l'autre.

%------------------------------------------------------------------
\partie{4. Le gradient spectral sur des images simples}{2:40 -- 4:00}{\texttt{gibbs.png}.}

Pour comprendre la méthode, j'ai utilisé une image dont on connaît tout : un carré clair sur fond sombre.

Sur un bord net, le gradient spectral oscille : 104, puis $-46$, 29, $-21$, et ça ne décroît qu'en
$1/k$. C'est l'effet de Gibbs : pour fabriquer un saut avec des sinusoïdes, la courbe interpolée doit
onduler. J'ai retrouvé exactement ces valeurs par le calcul : ce sont les restes de la série harmonique
alternée, en commençant par $\Delta \ln 2$. Dès que le bord est un peu flou, les oscillations
disparaissent.

Ces oscillations se compensent dans la reconstruction, mais apparaissent dès qu'on modifie le gradient.
Enfin, avec du bruit, le gradient spectral multiplie l'écart type par 1,8, contre 1,4 pour une différence
finie, comme le prévoit la théorie.

%------------------------------------------------------------------
\partie{5. Rehaussement de contraste et zones noires}{4:00 -- 5:30}{\texttt{contraste\_simple.png},
en insistant sur le profil (d).}

La première application de l'article est le rehaussement de contraste : on multiplie par $\alpha = 2{,}5$
le gradient des pixels plus sombres qu'un seuil $T$. Le contraste dans le sombre est bien multiplié par
2,5. Mais j'ai observé un défaut que l'article ne mentionne pas : une zone sombre entourée de clair
devient noire, elle passe même sous zéro.

L'explication est le bord. Le gradient spectral répartit le saut de part et d'autre du bord : la moitié
du côté sombre est sous le seuil, donc multipliée par $\alpha$. Le champ demande un saut supplémentaire
de $(\alpha - 1)$ fois la moitié du saut, et comme la moyenne de l'image est fixée, ce saut est réparti :
le clair monte et le sombre descend. Ce n'est pas seulement un effet de Gibbs : avec un bord doux, la
zone sombre descend quand même.

Pour éviter ces zones noires, j'ai proposé un étirement min-max : on décale l'image pour que son minimum
soit 0, puis on la multiplie pour que son maximum soit 255. Plus rien n'est écrêté, la zone sombre
redevient grise. Cet étirement a deux défauts : il dépend de deux pixels seulement, souvent des
dépassements de Gibbs, ce qui peut rendre toute l'image terne ; et en couleur, il crée des taches dans
les zones presque noires, que j'ai corrigées en ajoutant une petite constante au dénominateur.

%------------------------------------------------------------------
\partie{6. Comparaison avec l'article}{5:30 -- 6:30}{\texttt{article\_contraste.png} puis
\texttt{article\_collage.png}.}

J'ai ensuite appliqué mon code aux images de l'article. Le masque de la zone sombre coïncide avec le sien
à 97 à 99\,\%. Ses résultats publiés sont plus clairs que l'image $u$ brute, ce qui confirme qu'il
applique un étirement qu'il ne décrit pas. Avec mon étirement min-max, la zone sombre est éclaircie comme
dans l'article ; sur le jockey, le ciel est un peu plus terne, c'est le défaut dont je viens de parler.

Pour le collage et le mélange de gradients, mes résultats sont visuellement identiques aux siens : le
bateau cloné, le faon effacé, le texte sur le bois, la texture sur la texture.

%------------------------------------------------------------------
\partie{7. Collage : ce qui se passe hors du masque}{6:30 -- 7:00}{\texttt{section5.png} (courbe de
l'écart en fonction de la distance).}

Comme on résout sur toute l'image, sans condition au bord du masque, l'extérieur est un peu modifié :
jusqu'à 59 niveaux de gris au bord, encore 7 à 60 pixels. Il n'y a pas d'oscillation, mais un décalage
lent. C'est la limite que l'article reconnaît lui-même.

%------------------------------------------------------------------
\partie{8. Application à l'imagerie médicale}{7:00 -- 8:30}{\texttt{thorax\_contraste.png},
\texttt{nodules.png}, \texttt{effacement.png}.}

J'ai appliqué la méthode à des images médicales. Sur une radio du thorax, le seuil de 50 de l'article ne
sélectionne presque rien ; avec un seuil automatique qui couvre 25\,\% de l'image, on sélectionne les
poumons, et la trame vasculaire devient nettement plus visible.

Le collage permet d'insérer une lésion : j'ai collé une zone de poumon avec des nodules dans un poumon
sain. Avec le collage simple, les côtes disparaissent dans la zone ; avec le mélange de gradients, elles
sont conservées. Avec la même image comme source, on peut aussi effacer les marqueurs écrits sur la
radio, sans couture.

Enfin, la méthode s'étend directement aux volumes, avec une FFT en 3D, et il faut bien traiter le volume
en entier : coupe par coupe, chaque coupe reçoit son propre décalage et des stries apparaissent. Mais sur
une IRM, le rehaussement n'apporte rien : la zone sombre est surtout l'air et le liquide, et l'image
devient grise. Cette méthode n'est pas adaptée à tous les types d'images.

%------------------------------------------------------------------
\partie{9. Conclusion}{8:30 -- 9:20}{la liste des limites de la discussion, puis une image finale
(par exemple \texttt{thorax\_contraste.png}).}

Pour conclure : la méthode est exacte, simple et rapide, et elle se généralise sans difficulté à la 3D.
Mais mon analyse montre ses limites : le gradient spectral oscille près des bords nets, le collage modifie
un peu l'extérieur du masque, le rehaussement de contraste peut noircir les zones sombres et amplifie le
bruit autant que le signal. Mon étirement min-max évite les zones noires, mais dépend des pixels
extrêmes. Les suites naturelles seraient un étirement par percentiles, et le traitement des volumes à
pleine résolution. Merci.

%------------------------------------------------------------------
\bigskip
\hrule
\medskip
\textbf{Conseils pour l'enregistrement.}
\begin{itemize}
  \item Chronomètre-toi une première fois : si tu dépasses 9 min 30, raccourcis d'abord la partie 7
  (collage hors du masque), puis la partie 6.
  \item Pour chaque figure, montre du doigt (ou au curseur) ce dont tu parles : le profil sous zéro,
  les stries de l'IRM, les côtes dans la zone collée.
  \item Ne lis pas les formules en entier : montre l'équation 12 et dis « une simple division,
  coefficient par coefficient ».
  \item Si on te pose des questions sur l'IRM, voir \texttt{docs/explications/irm\_3d.pdf}.
\end{itemize}

\end{document}
